\section{Methodology}

\begin{table*}
    \centering
    \begin{tabular}{clcccc}
        \toprule
         Lock Name & Coherence & Access Modes & Starve-Free & Fair & RW-Safe \\
        \midrule
         Bakery Algorithm\cite{lamportNewSolutionDijkstras1974} & Software & Cached-SC, UC, NT & Yes & Yes & Yes\\
         Boulangerie\cite{mosesHoodBakeryAlgorithm2015} & Software & Cached-SC, UC, NT & Yes & Yes & Yes\\
         Lamport Fast\cite{lamportFastMutualExclusion1987} & Software & Cached-SC, UC, NT & No & No & No\\
         Elevator-LF-Array\cite{buhrHighcontentionMutualExclusion2018} & Software & Cached-SC, UC, NT & Yes & Yes & No\\
         Elevator-BL-Array\cite{buhrHighcontentionMutualExclusion2018} & Software & Cached-SC, UC, NT & Yes & Yes & Yes\\
         Elevator-LF-Tree\cite{buhrHighcontentionMutualExclusion2018} & Software & Cached-SC, UC, NT & Yes & Yes & No\\
         Elevator-BL-Tree\cite{buhrHighcontentionMutualExclusion2018} & Software & Cached-SC, UC, NT & Yes & Yes & Yes\\
         Peterson\cite{petersonMythsMutualExclusion1981} & Software & Cached-SC, UC, NT & No & Yes & No\\
         Knuth\cite{knuthAdditionalCommentsProblem1966} & Software & Cached-SC, UC, NT & No & Yes & No\\
         Spin\cite{andersonPerformanceSpinLock1990} & Hardware & Cached-HC & No & No & No\\
         Exp Spin\cite{andersonPerformanceSpinLock1990} & Hardware & Cached-HC & No & No & No\\
         Ticket\cite{mellor-crummeyAlgorithmsScalableSynchronization1991} & Hardware & Cached-HC & Yes & Yes & No\\
         MCS\cite{mellor-crummeyAlgorithmsScalableSynchronization1991} & Hardware & Cached-HC & Yes & Yes & No\\
         \bottomrule
    \end{tabular}
    \caption{Lock implementations and their applicable access modes. Hardware locks require cached access with hardware coherence (Cached-HC) for RMW atomics. Software locks are evaluated under cached software coherence (Cached-SC), uncached (UC), and non-temporal (NT) access modes.}
    \label{tab:Locks}
\end{table*}

\subsection{Benchmark Design}

We evaluate synchronization performance using microbenchmarks that measure lock throughput (successful acquisitions per second) under controlled contention.
We design three benchmark configurations to capture different contention regimes:

\paragraph{Max contention.}
All $N$ threads simultaneously compete to acquire the lock for a fixed time window.
When a thread acquires the lock, it executes a critical section and releases.
This represents the worst-case scenario for cache-line contention.

\paragraph{Moderate contention.}
$N$ threads each alternate between a non-critical computation phase (calibrated to 1\,$\mu$s) and a lock acquisition attempt.
This produces a mix of contended and uncontended acquisitions, representative of workloads where threads perform meaningful work between synchronization points.

\paragraph{Min contention.}
A single thread acquires and releases the lock repeatedly, measuring the inherent per-acquisition overhead as the lock is scaled to support $N$ threads.
This isolates the cost of lock metadata management (relevant for software locks with $O(N \log N)$ storage).

\paragraph{Critical section design.}
For max and moderate benchmarks, we evaluate two critical section sizes:
\begin{itemize}
    \item \textbf{Empty:} Increment a shared counter only. Isolates synchronization overhead.
    \item \textbf{Indexed:} Perform a hash-table lookup and update (64-byte key, 256-byte value) in CXL memory. Represents a realistic shared data structure operation where both the lock and the protected data reside in FAM.
\end{itemize}

\paragraph{Parameters.}
Thread count ranges from 1 to 64 in increments of 4.
Each configuration runs for 20 seconds (max/moderate) or 1 second (min), with a 1-second warmup.
Thread IDs are randomized between runs to mitigate bias from ID-dependent tie-breaking in some algorithms.
Each experiment is repeated 5 times; we report median throughput with error bars showing min/max.

\paragraph{Access mode configuration.}
For each lock and contention level, we measure performance under all applicable access modes (Table~\ref{tab:Locks}).
Cached-SC interposes explicit flush and invalidate operations around shared lock variable accesses.
After every store, we insert \texttt{clflushopt} followed by \texttt{sfence} to flush the modified cache line and order the writeback; before reads in polling loops, we insert \texttt{clflush} followed by \texttt{lfence} to invalidate the line and ensure the next load fetches from CXL memory.
These operations are implemented as compile-time macros (\texttt{FLUSH}/\texttt{INVALIDATE}) enabled via the \texttt{cached\_sc} build flag, producing zero overhead when disabled.
UC maps CXL memory uncacheable by opening the device file (e.g., \texttt{/dev/dax}) and calling \texttt{mmap(MAP\_SHARED)}; the kernel/driver configures the page table's PAT entry as UC.
Every load and store bypasses the host cache, so no explicit flushes or fences are needed on x86-64 (UC access is strongly ordered).
Cached-HC uses a standard writeback (WB) mapping with BISnp-managed coherence; RMW atomics are available.
NT access uses \texttt{movntdq}/\texttt{movntdqa} intrinsics for lock variable stores/loads, with an \texttt{sfence} after each store to flush the write-combining buffer; critical section data uses standard cached access.

\subsection{Mutex Implementations}

We select 13 mutex implementations spanning a range of design strategies (Table~\ref{tab:Locks}).

\paragraph{Why these locks?}
Our selection is guided by two criteria: (1)~covering the design space of contention management strategies (scanning, queueing, contention-avoidance), and (2)~including locks with different RW-safety properties relevant to CXL's ambiguous atomic load/store support.

\paragraph{Software locks.}
The \textbf{Bakery Algorithm}~\cite{lamportNewSolutionDijkstras1974} is Lamport's classical fair lock.
Each thread draws a ``ticket number'' by reading all other threads' numbers and choosing one larger than the maximum.
Threads enter the critical section in ticket order.
The algorithm is RW-safe: it tolerates non-atomic reads and writes by using a ``choosing'' flag to indicate when a thread is in the process of selecting its ticket.
This makes it directly applicable to CXL software-managed coherence without relying on atomic load/store guarantees.
\textbf{Boulangerie}~\cite{mosesHoodBakeryAlgorithm2015} optimizes Bakery by eliminating redundant reads of the choosing array when a thread can determine that no other thread is currently choosing.
\textbf{Lamport Fast}~\cite{lamportFastMutualExclusion1987} achieves the theoretical minimum of 5 memory operations in the uncontended fast path (2 stores, 3 loads) but is RW-unsafe and neither fair nor starvation-free.
\textbf{Peterson}~\cite{petersonMythsMutualExclusion1981} and \textbf{Knuth}~\cite{knuthAdditionalCommentsProblem1966} represent classical designs with extensive thread-state scanning.
Peterson's $N$-thread generalization uses $N-1$ levels of two-thread tournaments; Knuth's algorithm uses a circular scan of thread flags.
Both require reading $O(N)$ variables per acquisition attempt.

The \textbf{Elevator} series~\cite{buhrHighcontentionMutualExclusion2018} (LF-Array, BL-Array, LF-Tree, BL-Tree) are modern contention-avoiding designs.
Unlike classical software locks where all threads race to detect when the lock is free, Elevator locks have the unlock path explicitly wake the successor.
The successor search is performed either by a linear array scan (Array variants) or a binary tree traversal (Tree variants).
Each variant uses either Lamport-Fast (LF) or Burns-Lamport~\cite{burnsAllertonconf,hesselinkVerifyingSimplificationMutual2013} (BL) as the underlying mutex primitive for the initial lock acquisition.
The BL variants are RW-safe; the LF variants trade RW-safety for fewer memory operations.

\paragraph{Hardware locks.}
\textbf{Spin lock}~\cite{andersonPerformanceSpinLock1990} uses test-and-set (\texttt{lock xchg} on x86) with continuous polling on a single shared variable.
Every failed acquisition generates a RMW atomic on the contested cache line.
\textbf{Exp Spin} adds exponential back-off: after each failed test-and-set, the thread waits for an exponentially increasing number of cycles before retrying, reducing coherence traffic at the cost of increased acquisition latency.
\textbf{Ticket lock}~\cite{mellor-crummeyAlgorithmsScalableSynchronization1991} uses fetch-and-add (\texttt{lock xadd}) for fair ordering.
On acquire, a thread atomically increments the ``next ticket'' counter; it then polls the ``now serving'' counter with simple loads until its ticket is called.
This separates the RMW operation (one per acquisition) from the polling path (simple loads), reducing coherence traffic relative to spin locks.
\textbf{MCS lock}~\cite{mellor-crummeyAlgorithmsScalableSynchronization1991} queues per-thread lock objects using compare-and-swap (\texttt{lock cmpxchg}).
Each thread polls its own per-thread flag, so polling traffic is entirely local---no cache-line contention from concurrent pollers.
The predecessor sets the successor's flag during unlock, achieving $O(1)$ coherence traffic per acquisition regardless of thread count.

\paragraph{Memory ordering.}
All implementations use the minimum fences required for correctness under the x86-64 TSO memory model.
On x86-64, stores are not reordered with older stores, and loads are not reordered with older loads; thus most fence requirements reduce to \texttt{mfence} or \texttt{sfence} at store-release points.
We document fence placement for each lock implementation in our open-source benchmark~\cite{ramosEIRNfMutex_benchmark2025}.

\paragraph{Storage overhead.}
Software locks require $O(N)$ to $O(N \log N)$ storage where $N$ is the maximum thread count, as each thread needs dedicated state variables.
Hardware locks require $O(1)$ storage (a single cache line for spin/ticket, one cache line per thread for MCS).
At 64 threads, the largest software lock (Elevator-BL-Tree) requires 4.2\,KB; hardware locks require at most 4.1\,KB (MCS with 64-byte per-thread nodes).
This overhead is negligible relative to the protected data in any realistic application.
